\documentclass[11pt,a4paper]{article}
\usepackage[utf8]{inputenc}
\usepackage[T1]{fontenc}
\usepackage[french]{babel}
\usepackage{amsmath,amssymb}
\usepackage{booktabs}
\usepackage{geometry}
\usepackage{enumitem}
\usepackage{xcolor}
\usepackage{tcolorbox}
\usepackage{fancyhdr}

\geometry{margin=2cm}
\pagestyle{fancy}
\fancyhead[L]{\textbf{Reinforcement Learning -- Formulaire}}
\fancyhead[R]{EFREI 2024/2025}
\fancyfoot[C]{\thepage}

\newtcolorbox{formbox}[1][]{colback=blue!5,colframe=blue!50!black,title=#1,fonttitle=\bfseries}
\newtcolorbox{exbox}[1][]{colback=green!5,colframe=green!50!black,title=#1,fonttitle=\bfseries}
\newtcolorbox{tipbox}[1][]{colback=orange!5,colframe=orange!50!black,title=#1,fonttitle=\bfseries}

\newcommand{\E}{\mathbb{E}}

\begin{document}

\begin{center}
{\LARGE \textbf{Toutes les formules du RL}}\\[0.3cm]
{\large Avec les différentes écritures et quand les utiliser}
\end{center}

\vspace{0.5cm}

%% ============================================================
\section{Gain $G$ (3 écritures)}
%% ============================================================

\begin{formbox}[Forme 1 -- Somme longue (la plus intuitive)]
$$G_t = r_t + \gamma \, r_{t+1} + \gamma^2 \, r_{t+2} + \gamma^3 \, r_{t+3} + \cdots$$
\end{formbox}

\begin{formbox}[Forme 2 -- Somme compacte $\Sigma$ (la plus propre pour l'examen)]
$$\boxed{G_t = \sum_{k=0}^{\infty} \gamma^k \, r_{t+k}}$$
\end{formbox}

\begin{formbox}[Forme 3 -- Récursion (celle qui mène à Bellman)]
$$G_t = r_t + \gamma \, G_{t+1}$$
\end{formbox}

\begin{exbox}[Exemple numérique -- $\gamma = 0.9$, récompenses $[10, 5, 3, 1]$]
\begin{align*}
G_3 &= 1 \\
G_2 &= 3 + 0.9 \times 1 = 3.9 \\
G_1 &= 5 + 0.9 \times 3.9 = 8.51 \\
G_0 &= 10 + 0.9 \times 8.51 = \mathbf{17.66}
\end{align*}
\end{exbox}

%% ============================================================
\section{Value Function $V^\pi(s)$ (4 écritures)}
%% ============================================================

\begin{formbox}[Forme 1 -- Définition (pour Q3a de l'examen)]
$$\boxed{V^\pi(s) = \E\left[G_t \;\middle|\; S_t = s\right]}$$
\textit{``La moyenne du gain en partant de l'état $s$, en suivant la politique $\pi$.''}
\end{formbox}

\begin{formbox}[Forme 2 -- Somme développée]
$$V^\pi(s) = \E\left[\sum_{k=0}^{\infty} \gamma^k \, r_{t+k} \;\middle|\; S_t = s\right]$$
C'est la forme 1 où l'on remplace $G_t$ par sa définition.
\end{formbox}

\begin{formbox}[Forme 3 -- Bellman Expectation (pour Q15)]
$$\boxed{V^\pi(s) = \E\left[r_t + \gamma \, V^\pi(S_{t+1}) \;\middle|\; S_t = s\right]}$$
\textit{``Récompense immédiate $+\;\gamma \times$ valeur de l'état suivant.''}
\end{formbox}

\begin{formbox}[Forme 4 -- Bellman développée (avec les probas)]
$$V^\pi(s) = \sum_{a} \pi(a|s) \left[\sum_{r} r \, p(r|s,a) \;+\; \gamma \sum_{s'} p(s'|s,a) \, V^\pi(s')\right]$$
C'est la forme 3 où l'espérance est écrite explicitement.
\end{formbox}

\begin{formbox}[Bellman Optimality (pour la politique optimale)]
$$\boxed{V^*(s) = \max_a \; \E\left[r_t + \gamma \, V^*(S_{t+1}) \;\middle|\; S_t = s, \, A_t = a\right]}$$
\end{formbox}

%% ============================================================
\section{Action-Value Function $Q^\pi(s,a)$ (4 écritures)}
%% ============================================================

\begin{formbox}[Forme 1 -- Définition (pour Q3b de l'examen)]
$$\boxed{Q^\pi(s,a) = \E\left[G_t \;\middle|\; S_t = s, \, A_t = a\right]}$$
\textit{``La moyenne du gain en partant de $s$, en prenant l'action $a$, puis en suivant $\pi$.''}
\end{formbox}

\begin{formbox}[Forme 2 -- Via $V$ (one-step lookahead)]
$$Q^\pi(s,a) = \E\left[r_t + \gamma \, V^\pi(S_{t+1}) \;\middle|\; S_t = s, \, A_t = a\right]$$
\textit{``Récompense $+\;\gamma \times V$ de l'état d'après.''}
\end{formbox}

\begin{formbox}[Forme 3 -- Bellman pour $Q$ (base de SARSA)]
$$\boxed{Q^\pi(s,a) = \E\left[r_t + \gamma \, Q^\pi(S_{t+1}, A_{t+1}) \;\middle|\; S_t = s, \, A_t = a\right]}$$
\textit{``Récompense $+\;\gamma \times Q$ de l'état-action d'après.''}
\end{formbox}

\begin{formbox}[Forme 4 -- Bellman Optimality pour $Q^*$ (base de Q-Learning)]
$$\boxed{Q^*(s,a) = \E\left[r_t + \gamma \, \max_{a'} Q^*(S_{t+1}, a') \;\middle|\; S_t = s, \, A_t = a\right]}$$
\textit{``Récompense $+\;\gamma \times$ le meilleur $Q$ possible après.''}
\end{formbox}

%% ============================================================
\section{Lien entre $V$, $Q$ et $G$}
%% ============================================================

\begin{formbox}[Relations fondamentales]
\begin{align}
V^\pi(s) &= \sum_a \pi(a|s) \; Q^\pi(s,a) && \text{$V$ = moyenne de $Q$ sur les actions} \\[6pt]
Q^\pi(s,a) &= \E\left[r + \gamma \, V^\pi(s') \;\middle|\; s, a\right] && \text{$Q$ = récompense + $\gamma \, V$ suivant} \\[6pt]
\pi^*(s) &= \arg\max_a \; Q^*(s,a) && \text{Politique optimale depuis $Q^*$} \\[6pt]
V^*(s) &= \max_a \; Q^*(s,a) && \text{$V^*$ = le meilleur $Q^*$}
\end{align}
\end{formbox}

\begin{exbox}[Analogie -- Carrefour routier]
\begin{itemize}[nosep]
\item $G$ = le temps total d'UN trajet précis
\item $V(s)$ = ``ce carrefour est bien situé en moyenne'' (score du lieu)
\item $Q(s, a)$ = ``ce carrefour est bien SI tu vas à gauche'' (score du lieu + direction)
\item $V(s) = 0.5 \times Q(s, \text{gauche}) + 0.5 \times Q(s, \text{droite})$ si politique 50/50
\end{itemize}
\end{exbox}

%% ============================================================
\section{Règles de mise à jour des algorithmes}
%% ============================================================

\begin{formbox}[Monte Carlo (MC) -- pas de bootstrapping]
$$V(S_t) \;\leftarrow\; V(S_t) + \alpha \Big[\underbrace{G_t}_{\text{retour complet}} - V(S_t)\Big]$$
\end{formbox}

\begin{formbox}[TD(0) -- avec bootstrapping]
$$V(S_t) \;\leftarrow\; V(S_t) + \alpha \Big[\underbrace{r_t + \gamma \, V(S_{t+1})}_{\text{TD target}} - V(S_t)\Big]$$
\end{formbox}

\begin{formbox}[SARSA -- On-Policy TD Control]
$$\boxed{Q(S_t, A_t) \;\leftarrow\; Q(S_t, A_t) + \alpha \Big[r_t + \gamma \, Q(S_{t+1}, \textcolor{red}{A_{t+1}}) - Q(S_t, A_t)\Big]}$$
$A_{t+1}$ = action \textbf{réellement prise} ($\varepsilon$-greedy) $\;\Rightarrow\;$ \textbf{on-policy}
\end{formbox}

\begin{formbox}[Q-Learning -- Off-Policy TD Control]
$$\boxed{Q(S_t, A_t) \;\leftarrow\; Q(S_t, A_t) + \alpha \Big[r_t + \gamma \, \textcolor{red}{\max_{a'}} \, Q(S_{t+1}, a') - Q(S_t, A_t)\Big]}$$
$\max_{a'}$ = meilleure action \textbf{théorique} $\;\Rightarrow\;$ \textbf{off-policy}
\end{formbox}

\begin{formbox}[REINFORCE -- Policy Gradient (Monte Carlo)]
$$\theta \;\leftarrow\; \theta + \alpha \, \nabla_\theta \log \pi_\theta(A_t | S_t) \; G_t$$
\end{formbox}

\begin{formbox}[Policy Gradient Theorem]
$$\boxed{\nabla_\theta J(\theta) = \E\left[\sum_t \nabla_\theta \log \pi_\theta(A_t | S_t) \; G_t\right]}$$
Permet d'estimer le gradient \textbf{sans connaître le modèle} de l'environnement.
\end{formbox}

%% ============================================================
\section{Tableau récapitulatif -- Quelle formule pour quelle question}
%% ============================================================

\begin{center}
\renewcommand{\arraystretch}{1.4}
\begin{tabular}{lll}
\toprule
\textbf{Question examen} & \textbf{Ce qu'on écrit} & \textbf{Forme} \\
\midrule
Q3a -- Définition de $V$ & $V^\pi(s) = \E[G_t \mid S_t = s]$ & Définition \\
Q3b -- Définition de $Q$ & $Q^\pi(s,a) = \E[G_t \mid S_t = s, A_t = a]$ & Définition \\
Q3c -- Définition de $G$ & $G_t = \sum_{k=0}^{\infty} \gamma^k r_{t+k}$ & Somme $\Sigma$ \\
Q15 -- Équation de Bellman & $V^\pi(s) = \E[r + \gamma V^\pi(s') \mid s]$ & Bellman Expect. \\
Q16 -- Calcul de $V$ & $V(s) = R(s) + \gamma \sum_{s'} P(s'|s) V(s')$ & Bellman développée \\
Q21 -- Q-Learning & $Q(s,a) \leftarrow Q + \alpha[r + \gamma \max Q(s',a') - Q]$ & Mise à jour \\
Q29 -- Policy Gradient & $\nabla J = \E[\nabla \log \pi \cdot G_t]$ & PG Theorem \\
\bottomrule
\end{tabular}
\end{center}

\begin{tipbox}[Conseil examen]
Pour les questions de \textbf{définition} (Q3a, Q3b, Q3c) : utilise les formes avec $\E[G|\cdots]$, c'est les plus courtes et claires.\\
Garde les formes \textbf{Bellman} pour la question 15 et les calculs (Q16).\\
Garde les formes de \textbf{mise à jour} pour identifier les algorithmes (Q21).
\end{tipbox}

\end{document}
